\subsection*{Note (6)}
The sixth note, beginning on p.~55 and ending on p.~56, qualifies the
tame-topology discussion attached to \S5, p.~31. Grothendieck identifies
the isomorphism classes of compact tame spaces with those of the
piecewise linear theory, while recalling the stability requirements that
motivate his earlier objections to the latter as a complete context for
the program. He presents this as a rehabilitation, in the intended
setting, of the ``Hauptvermutung''. The important change of emphasis
comes on p.~56: the polemic against wild phenomena as obstacles in the
chosen foundations is accompanied by an explicit recognition of their
reason for being and their own beauty. The proposed foundations are to
place them as ``cas limites'' rather than make them disappear. This
reassessment belongs beside the construction and comparison requirements
of \S5, pp.~27--33, and the treatment of full homotopy types in \S6,
pp.~38--39. The note states the intended comparison and its interpretation;
the comparison's proof is not developed here.
(\href{https://grothendieck.umontpellier.fr/web/viewer.html?file=../119.pdf\#page=58}{Copy 1, pp.~55--56, beginning at site 57};
\href{https://grothendieck.umontpellier.fr/web/viewer.html?file=../119.pdf\#page=59}{closing qualification on p.~56, site 58};
\href{https://grothendieck.umontpellier.fr/web/viewer.html?file=../119.pdf\#page=34}{call in \S5, p.~31, site 33}.)
